\documentclass[conference]{IEEEtran}

\usepackage{multirow}
\usepackage{float}
\usepackage{cite}
\usepackage{caption}
\usepackage{subcaption}
\usepackage{amsmath,amssymb,amsfonts,amstext}
% \usepackage{algorithmic}
\usepackage[dvipdfmx]{graphicx}
\usepackage{textcomp}
\usepackage{xcolor}
\usepackage{colortbl}
\usepackage[T1]{fontenc}
\usepackage{mdframed}
\usepackage{algorithm, algpseudocode} % texlive-science
\usepackage{newtxtext}
\usepackage[normalem]{ulem}
\usepackage{nameref}
\usepackage{tikz}
\usetikzlibrary{arrows.meta, shapes, positioning, calc, fit}
\usepackage{quantikz}
%\newcommand*\rtofstarg{\tikz[baseline=(char.base)]{
            \node[shape=circle,double,draw,inner sep=1pt,minimum width=16pt] (char) {S};}}
\newcommand*\rtofsdgtarg{\tikz[baseline=(char.base)]{
            \node[shape=circle,double,draw,inner sep=1pt] (char) {S$^{\dagger}$};}}
\newcommand*\rtoftarg{\tikz[baseline=(char.base)]{
            \node[shape=circle,draw,inner sep=-0.5pt,fill=none] (char) {$\oplus$};}}
\newcommand*\rtofsctrl{\tikz[baseline=(char.base)]{
            \node[shape=circle,double,draw,inner sep=1pt,fill=black, minimum width=3pt] (char) {};}}
\newcommand*\rtofsctrlo{\tikz[baseline=(char.base)]{
            \node[shape=circle,double,draw,inner sep=1pt,fill=white, minimum width=3pt] (char) {};}}
%\newcommand*\rtofsctrlo{\tikz[baseline=(char.base)]{
%            \node[shape=circle,double,draw,inner sep=1pt,fill=white, minimum width=3pt] (char) {};}}
\newcommand*\dashedwire{\tikz{
            \draw[dashed] (0,0) -- (0.5,0); }}

\newcommand*\ctrlrtofs[1]{\push{\rtofsctrl}\qwx[#1]\qw}
\newcommand*\ctrlortofs[1]{\push{\rtofsctrlo}\qwx[#1]\qw}
\newcommand*\targrtofs{\push{\rtofstarg}\qw}
\newcommand*\targrtof{\push{\rtoftarg}\qw}
% \newcommand*\qcbox[1,2,3,4]{\gategroup[4,steps=2,style={dashed,rounded corners,fill=blue!20, inner xsep=2pt},background,label style={label position=below,anchor=north,yshift=-0.2cm}]{Compute}}


\makeatletter
\newcommand{\sublabel}[1]{\protected@edef\@currentlabel{\thefigure(\thesubfigure)}\label{#1}}
\makeatother
\captionsetup{subrefformat=parens}
\algrenewcommand\algorithmicindent{0.5em}
\newtheorem{definition}{Definition}
\newtheorem{example}{Example}
\newcommand{\thickhline}{%
    \noalign {\ifnum 0=`}\fi \hrule height 1pt
    \futurelet \reserved@a \@xhline
}
%\def\BibTeX{{\rm B\kern-.05em{\sc i\kern-.025em b}\kern-.08em
%    T\kern-.1667em\lower.7ex\hbox{E}\kern-.125emX}}
\newcommand{\figcaption}[1]{\def\@captype{figure}\caption{#1}}
\newcommand{\tblcaption}[1]{\def\@captype{table}\caption{#1}}

\newcommand{\Del}[1]{\textcolor{red}{\sout{#1}}}
\newcommand{\Add}[1]{\textcolor{blue}{\emph{#1}}}

\newcount\NGCcolor
\NGCcolor=1
\newcount\NGCcolortwo
\NGCcolortwo=1
%\NGCcolor=0
\ifnum\NGCcolortwo<1
  \newcommand{\blueoutt}[1]{\textcolor{blue}{#1}}  
\else
  \newcommand{\blueoutt}[1]{#1}
\fi
  
\ifnum\NGCcolor<1
  \newcommand{\redout}[1]{\textcolor{red}{#1}}
  \newcommand{\blueout}[1]{\textcolor{blue}{#1}}
  \newcommand{\greenout}[1]{\textcolor{green}{#1}}
  \newcommand{\cyannout}[1]{\textcolor{cyan}{#1}}
  \newcommand{\magentaout}[1]{\textcolor{magenta}{#1}}
  \newcommand{\yellowout}[1]{\textcolor{yellow}{\bf #1}}
  \newcommand{\algblueout}[1]{\leavevmode\color{blue}{#1}}
\else
  \newcommand{\redout}[1]{#1}
  \newcommand{\blueout}[1]{#1}
  \newcommand{\greenout}[1]{#1}
  \newcommand{\cyannout}[1]{#1}
  \newcommand{\magentaout}[1]{#1}
  \newcommand{\yellowout}[1]{#1}
  \newcommand{\algblueout}[1]{\leavevmode {#1}}
\fi

% refs for section
\newcommand{\secref}[1]{Sec.~\ref{#1}}
\newcommand{\figref}[1]{Fig.~\ref{#1}}
\newcommand{\tabref}[1]{Table~\ref{#1}}
\newcommand{\eqnref}[1]{(\ref{#1})}
%% \newcommand*\rtofstarg{\tikz[baseline=(char.base)]{
%%             \node[shape=circle,double,draw,inner sep=1pt,minimum width=16pt] (char) {S};}}
%% \newcommand*\rtofsdgtarg{\tikz[baseline=(char.base)]{
%%             \node[shape=circle,double,draw,inner sep=1pt] (char) {S$^{\dagger}$};}}
%% \newcommand*\rtoftarg{\tikz[baseline=(char.base)]{
%%             \node[shape=circle,double,draw,inner sep=1pt] (char) {\phantom{S$^{\dagger}$}};}}
%% \newcommand*\rtofsctrl{\tikz[baseline=(char.base)]{
%%             \node[shape=circle,double,draw,inner sep=1pt,fill=black, minimum width=3pt] (char) {};}}
%% \newcommand*\rtofsctrlo{\tikz[baseline=(char.base)]{
%%             \node[shape=circle,double,draw,inner sep=1pt,fill=white, minimum width=3pt] (char) {};}}

%% \newcommand*\ctrlrtofs[1]{\push{\rtofsctrl}\qwx[#1]\qw}
%% \newcommand*\ctrlortofs[1]{\push{\rtofsctrlo}\qwx[#1]\qw}
%% \newcommand*\targrtofs{\push{\rtofstarg}\qw}
%% \newcommand*\targrtof{\push{\rtoftarg}\qw}

\begin{document}

\scalebox{0.6}{
  \begin{quantikz}[row sep=0.2em, column sep=0.2em]
\lstick{$x_1$}     &                   &                      & \gate[3]{F}          &                   &                      & \qw               \\
\lstick{$x_2$}     & \gate[2]{A}       &                      &                      &                   &                      & \qw               \\
\lstick{$x_3$}     & \ghost{A}\ctrl{2} & \gate[1]{C}\ctrl{3}  &                      & \gate[2]{B}       & \gate[3]{D}          & \qw                 \\
\lstick{$x_4$}     &                   &                      &                      & \ghost{B}\ctrl{2} & \ghost{D}\ctrl{2}\qw & \qw               \\
\lstick{$\ket{0}$} &                   & \targ{}              &                      &                   & \ghost{D}\ctrl{2}    & \qw               \\
\lstick{$\ket{0}$} &                   & \targX{}             &                      &                   & \targ{}              & \qw               \\
\lstick{$\ket{0}$} &                   &                      &                      &                   &                      & \qw               \\
\lstick{$\ket{0}$} &                   &                      &                      &                   &                      & \qw               \\
\end{quantikz}

%% \begin{quantikz}[row sep=0.3cm, column sep=0.55cm]
%% \qw              & \qw             & \qw        & \qw             & \qw              & \gate[2]{F_i}       & \qw                    & \qw       & \qw             & \qw      & \qw        & \qw              & \gate[2]{F_i}        & \qw                     & \qw        \\
%% \qw              & \qw             & \qw        & \qw             & \qw              & \ghost{F_i}         & \qw                    & \qw       & \qw             & \qw      & \qw        & \qw              & \ghost{F_i}          & \qw                     & \qw        \\
%% \qw              & \qw             & \qw        & \qw             & \ctrl{2}         & \qw                 & \ctrl{2}               & \qw       & \qw             & \qw      & \qw        & \ctrl{2}         & \qw                 & \ctrl{2}                & \qw         \\
%% \qw              & \qw             & \octrl{2}  & \qw             & \qw              & \qw                 & \qw                    & \qw       & \qw             & \qw      & \qw        & \qw              & \qw                 & \qw                     & \qw         \\
%% \lstick{$a_1$}   & \qw             & \qw        & \qw             & \gate{i\omega Z} & \targ{}             & \gate{i\omega Z^\dagger} & \qw       & \octrl{1}       & \qw      & \qw        & \gate{i\omega Z} & \targ{}             & \gate{i\omega Z^\dagger}  & \qw          \\
%% \lstick{$a_2$}   & \qw             & \gate{iZ}  & \gate{H}        & \ctrl{-1}        & \qw                 & \ctrl{-1}              & \gate{H}  & \gate{iZ^\dagger} & \qw      & \gate{H}   & \ctrl{-1}        & \qw                 & \ctrl{-1}               & \gate{H}     \\
%% \lstick{}        & \gate{H}        & \ctrl{-1}  & \qw             & \qw              & \qw                 & \qw                    & \qw       & \ctrl{-1}       & \gate{H} & \qw        & \qw              & \qw                 & \qw                     & \qw          \\
%%\end{quantikz}


  %\input{img/fig-func-step}
}

\end{document}
